\section{Reduction semantics and structural congruence}\label{sec:red}

\subsection{Structural congruence}

Processes are identified up to $\alpha$-conversion of restricted channels and of
bound variables. The structural congruence $\equiv$ is the least equivalence
relation containing the rules of \cref{tab:cgr} and closed under all
contexts: conditionals, parallel composition, restriction, recursion,
prefixes and sums. Congruent processes have the same tags.

\begin{table}[h]
\centering
\small
$\begin{array}{l@{\quad}rcl@{\quad}l}
\textsc{ParNil} & P \parallel \mathbf{0} &\equiv& P \\
\textsc{ParComm} & P \parallel Q &\equiv& Q \parallel P \\
\textsc{ParAssoc} & P \parallel (Q \parallel R) &\equiv& (P \parallel Q) \parallel R \\[2pt]
\textsc{SumNil} & K + \mathbf{0}_{\tilde e} &\equiv& K \\
\textsc{SumComm} & K + K' &\equiv& K' + K \\
\textsc{SumAssoc} & K + (K' + K'') &\equiv& (K + K') + K'' \\[2pt]
\textsc{IteT} & \ite{e}{P}{Q} &\equiv& P & e \Downarrow \mathit{tt} \\
\textsc{IteF} & \ite{e}{P}{Q} &\equiv& Q & e \Downarrow \mathit{ff} \\
\textsc{ValExpr} & P &\equiv& P[v/e] & e \Downarrow v \\[2pt]
\textsc{ResNil} & \mathbf{0}_{\tilde e} \setminus c &\equiv& \mathbf{0}_{\tilde e} \\
\textsc{ResOrd} & (P \setminus c) \setminus d &\equiv& (P \setminus d) \setminus c \\
\textsc{ResPar} & (P \parallel Q) \setminus c &\equiv& P \parallel (Q \setminus c) & c \notin \fc(P)
\end{array}$
\caption{Structural congruence. In \textsc{ValExpr}, $e$ is an expression of
an output, of a conditional, of a deadlock $\mathbf{0}_{\tilde e}$ or of a
quantum prefix, and $P[v/e]$ replaces it by its value. $\fc(P)$ is the set of
free channels of $P$.}
\label{tab:cgr}
\end{table}

\begin{remark}\label{rem:debruijn}
Scope extrusion \textsc{ResPar} needs the $\alpha$-conversion of the restricted
channel: $(R \setminus d) \parallel Q$ with $d$ free in $Q$ must first rename
$d$. As in VACCS, the mechanisation uses de Bruijn indices for restricted
channels and for the variables bound by inputs and measurements:
\textsc{ResPar} reads $(\nu P) \parallel Q \equiv \nu (P \parallel Q^{+})$,
where $Q^{+}$ shifts the indices of $Q$, \textsc{ResOrd} exchanges the two
innermost indices, and the substitution $P[Q/X]$ shifts $Q$ under each binder,
so that $\equiv$ is closed under recursion.
\end{remark}

\subsection{Reduction semantics}

The reduction semantics of configurations, written
$\conf{\rho}{P} \red_{s} \conf{\rho'}{P'}$ and indexed by the choice $s$ that
selects the redex, is the least relation closed under the rules of
\cref{fig:sts}. It follows the STS of VACCS: the first six rules are the
redexes, the last three close them under parallel composition, restriction
and structural congruence.

\begin{figure}[h]
\small
\begin{mathpar}
\inferrule*[left=Com]
  {e \Downarrow v}
  {\conf{\rho}{\tagged{t}{c!e} \parallel (\tagged{t'}{c?x}.P + K)} \red_{\{t,t'\}}
   \conf{\rho}{\mathbf{0} \parallel P[v/x]}}

\inferrule*[left=Tau]
  { }
  {\conf{\rho}{\tagged{t}{\tau}.P + K} \red_{t} \conf{\rho}{P}}

\inferrule*[left=Op]
  {U : \mathsf{Op}(|\tilde q|)}
  {\conf{\rho}{\tagged{t}{U(\tilde q)}.P + K} \red_{t} \conf{\upd{U}}{P}}

\inferrule*[left=Init]
  { }
  {\conf{\rho}{\tagged{t}{\initp}.P + K} \red_{t} \conf{\inirho}{P}}

\inferrule*[left=Meas]
  {\M : \mathsf{Meas}(|\tilde q|) \\ 0 \le m < |\M|}
  {\conf{\rho}{\tagged{t}{\M(\tilde q \rhd x)}.P + K} \red_{t} \conf{\br}{P[m/x]}}

\inferrule*[left=Rec]
  { }
  {\conf{\rho}{\tagged{t}{\rec{X}{P}}} \red_{t}
   \conf{\rho}{P[\tagged{t}{\rec{X}{P}}/X]}}

\inferrule*[left=Par]
  {\conf{\rho}{P} \red_{s} \conf{\rho'}{P'}}
  {\conf{\rho}{P \parallel Q} \red_{s} \conf{\rho'}{P' \parallel Q}}

\inferrule*[left=Res]
  {\conf{\rho}{P} \red_{s} \conf{\rho'}{P'}}
  {\conf{\rho}{P \setminus c} \red_{s} \conf{\rho'}{P' \setminus c}}

\inferrule*[left=Cong]
  {P \equiv P_1 \\ \conf{\rho}{P_1} \red_{s} \conf{\rho'}{Q_1} \\ Q_1 \equiv Q}
  {\conf{\rho}{P} \red_{s} \conf{\rho'}{Q}}
\end{mathpar}
\caption{Reduction semantics of configurations. A measurement reduces to one
unnormalised branch for each outcome $m$; all branches carry the same choice
$t$.}
\label{fig:sts}
\end{figure}
